Abschlussarbeit ist wie ein Rezept
-> jemand anderes sollte zum gleichen Ergebnis kommen
-> jemand der die Arbeit liest, soll sie mit der Anleitung reproduzieren können
wird beschrieben:
- alles was ein KOSI/KITS des gleichen Jahrgangs nicht wissen kann
  -> außerhalb des Unistoffs
  -> nur beschreiben, was man neu entworfen hat
keine Detaillierte Beschreibung von:
- alles andere / generelle entfällt oder nur kurz erwähnen, Quelle dran, fertig
- keine Firmenbeschreibung - Firma ist Mittel zum Zweck = Experimentierumgebung

Thesis-Text
===========
- keine Fragewörter in nicht-Fragen verwenden - nicht: "wie oben gezeigt"
- Denglisch erlaubt
  -> im üblichen Sprachgebrauch - nur Eigennamen, keine Verben
  -> Informatik eh englisch, aber konsistent bleiben, nicht hin und her wechseln
- Zitieren: von LateX fine - Buchstaben und Zahlen
- Inspiration für Einleitung/Abstract/Fazit \url{https://www.researchgate.net/publication/290368471_Der_Benutzungsfragebogen_ISONORM_924110_Ergebnisse_zur_Reliabilitat_und_Validitat}

Fragestellung
=============
- Warum, Wie und Was beantworten
- notfalls eine Unterüberschrift benutzen (= gängige Methode zur genaueren Beschreibung)
- Ist mein Thema mit der FF vollständig abgedeckt?
- Ist die Frage zu allgemein und sagt mehr aus, als das was ich mache?
  -> potentiell in 2 Subfragen aufsplitten
- Domänenfixiert: Ich habe diese Arbeit in dieser Firma umgesetzt
  -> ist dann aber auch nur dort möglich

Quellen
=======
- Tool finden für Quellengestützte Notizen (Obsidian statt OneNote)
- Quelle: ist gut gearbeitet und ist es wissenschaftlich gerated worden
  -> Master/Doctorarbeiten sind auch Quellen (BA nicht empfohlen)

Empfehlung für die Endphase:
-	Formatierungshilfe – gerne abnehmen lassen (Bildunterschriten, Kommas, Einrückungen etc.)
